\subsection{ABC Notation}\label{subsec:abc-notation}

\textbf{ABC notation} is a text-based music notation system designed to be comprehensible by both humans and computers,
allowing musical scores to be represented using plain text characters such as letters, digits, and
punctuation~\cite{Walshaw_2011}.

It was developed as a lightweight alternative to traditional staff notation, enabling easy storage, sharing, and
processing of musical compositions in digital form.
The system encodes music declaratively, combining structured metadata with symbolic note representations.
An ABC file typically consists of a \textit{header}, which defines musical properties such as title, meter,
and key, and a \textit{tune body}, which contains the actual sequence of musical events~\cite{Walshaw_2011}.

%% ---------------------------------------------------------------------------------------------------------------------

\subsubsection{Key Features}

A central component of \textbf{ABC notation} is its use of \textit{information fields}, which provide structured
metadata about a musical piece.
These fields, such as \texttt{T:} for title, \texttt{M:} for meter, \texttt{L:} for unit note length, and \texttt{K:}
for key, allow the representation of contextual musical information alongside the score~\cite{Walshaw_2011}.

Musical notes are represented using letters in the range \texttt{A–G}, with octave modifiers applied through commas and
apostrophes.
Durations are expressed relative to a base unit length and can be modified using numeric multipliers or fractional
notation~\cite{Walshaw_2011}.

The language supports a variety of musical constructs, including chords, rests, tuplets, slurs, and ties.
Rhythmic variations can be expressed using constructs such as broken rhythm markers (\texttt{>} and \texttt{<}),
enabling compact encoding of common rhythmic patterns~\cite{Walshaw_2011}.
ABC notation also provides mechanisms for structuring compositions through bar lines, repeats, and multiple voices.
The \texttt{V:} field enables multi-voice compositions, while additional directives allow for annotations, lyrics, and
formatting instructions~\cite{Walshaw_2011}.

Furthermore, ABC notation is widely supported by external tools capable of converting textual representations into
typeset sheet music or MIDI playback, making it a versatile format for both documentation and performance.

%% ---------------------------------------------------------------------------------------------------------------------

\subsubsection{Limitations}

Despite its expressiveness as a notation system, \textbf{ABC notation} is not designed as a programming language and
lacks support for higher-level abstractions.
It does not provide constructs for modular composition, reusable patterns, or hierarchical structuring beyond basic
mechanisms such as macros and repeated sections.

Additionally, the language does not explicitly model time as a first-class concept, instead relying on implicit ordering
of notes and durations, which can make it difficult to represent complex temporal relationships.
The system also conflates multiple concerns, including musical representation, metadata, and layout directives, which
may reduce clarity when attempting to use it as a foundation for compositional systems.

Finally, ABC notation does not include capabilities for sound synthesis or real-time interaction, relying instead on
external tools for playback and rendering.

%% ---------------------------------------------------------------------------------------------------------------------

\subsubsection{Relevance to Our DSL}

\textbf{ABC notation} demonstrates how musical structures can be mapped into a clear and human-readable textual format,
providing valuable insight into the design of accessible music-oriented languages.
Its use of structured metadata and symbolic representations offers a strong foundation for expressing musical ideas in a
compact and interpretable way.

However, its limitations in abstraction, modularity, and temporal modeling highlight key areas addressed by the DSL
proposed in this thesis, which aims to provide a more expressive, structured, and executable approach to music
composition.
